% Problem record op_8cd24cb9148d9dab. This TeX file is derived from
% database/problems_json/op_8cd24cb9148d9dab.json by scripts/sync-tex.mjs; edit the JSON record
% and rerun the script rather than editing this file.
\section{Exact CNOT count for arbitrary three-qubit unitaries}
\label{sec:op_8cd24cb9148d9dab}

\paragraph{Problem.}
What is the smallest integer $k$ such that every three-qubit unitary can be implemented exactly, up to a global phase, by a circuit with at most $k$ CNOT gates?

The circuit acts on exactly the three data qubits. Arbitrary one-qubit unitaries are free, and any ordered pair of distinct qubits can be the control and target of a CNOT. Ancillas and measurements are excluded. The full circuit, including the sequence of control and target pairs, may depend on the target unitary.

\subsection*{Status}

Unsolved

\subsection*{Source}

Krol and Al-Ars list the remaining gap between the theoretical lower bound and their construction in Sec.~1, Table~1 \sourcecite{ref:d11-block-zxz}{KA24}.

\subsection*{Progress}

\begin{itemize}
  \item Shende, Bullock, and Markov give a parameter-counting lower bound of $\lceil(4^n-3n-1)/4\rceil$ CNOTs for general $n$-qubit unitaries. For three qubits, the unitary up to phase has $63$ real parameters and a circuit with $k$ CNOTs has at most $9+4k$, so $k\geq 14$ \sourcecite{ref:d11-logic-synthesis}{SBM06}.

  \item The block-ZXZ decomposition in Krol and Al-Ars uses $19$ CNOTs for arbitrary three-qubit targets. See Sec.~1, Table~1, and Sec.~5.3 \sourcecite{ref:d11-block-zxz}{KA24}.
\end{itemize}

\subsection*{References}

\begin{enumerate}[label={},leftmargin=4.8em,labelsep=0.5em]
  \footnotesize
  \item[\textup{[KA24]}]\label{ref:d11-block-zxz}
  A. M. Krol and Z. Al-Ars, \emph{Beyond Quantum Shannon: Circuit Construction for General n-Qubit Gates Based on Block ZXZ-Decomposition} (2024). \href{https://arxiv.org/abs/2403.13692v2}{arXiv:2403.13692v2}.

  \item[\textup{[SBM06]}]\label{ref:d11-logic-synthesis}
  V. V. Shende, S. S. Bullock, and I. L. Markov, \emph{Synthesis of Quantum Logic Circuits} (2006). \href{https://arxiv.org/abs/quant-ph/0406176v5}{arXiv:quant-ph/0406176v5}.

  \item[\textup{[Paw25]}]\label{ref:d11-triality}
  B. Pawlowski, \emph{Real 3-qubit gate decompositions via triality}. \href{https://arxiv.org/abs/2512.18049v2}{arXiv:2512.18049v2}.
\end{enumerate}

\subsection*{Comment}

A matching upper bound requires an exact circuit for every three-qubit unitary. Pawlowski gives a $14$-CNOT construction for $\mathrm{SO}(8)$, the determinant-one real orthogonal targets, in Theorem~5.4 \sourcecite{ref:d11-triality}{Paw25}. The catalog and literature were checked through November 26, 2026.

\subsection*{Field}

Quantum Computation

\subsection*{Topic}

Quantum circuit complexity

\subsection*{ID}

\texttt{op\_8cd24cb9148d9dab}
